%% Elevark fra mattelab.no – laget automatisk av bygg.py. Rediger gjerne til eget bruk.
%% Kompiler med lualatex.
\documentclass[11pt,a4paper]{article}
\usepackage[norsk]{babel}
\usepackage[a4paper,margin=2.2cm,top=2cm,bottom=2.4cm]{geometry}
\usepackage{fontspec}
\usepackage{unicode-math}
\setmainfont{TeX Gyre Pagella}
\setsansfont{TeX Gyre Heros}[Scale=MatchLowercase]
\setmathfont{TeX Gyre Pagella Math}
\usepackage{xcolor}
\definecolor{blekk}{HTML}{16213A}
\definecolor{dempet}{HTML}{566077}
\definecolor{rute}{HTML}{CFC9B8}
\definecolor{m1}{HTML}{2F4AB0}
\definecolor{m2}{HTML}{99266D}
\usepackage[most]{tcolorbox}
\usepackage{qrcode}
\usepackage{needspace}
\usepackage{fancyhdr}
\usepackage[hidelinks]{hyperref}
\pdfvariable suppressoptionalinfo \numexpr 1+2+4+8+16+32+64+128+256+512\relax
\setlength{\parindent}{0pt}
\setlength{\parskip}{0.4em}
\pagestyle{fancy}
\fancyhf{}
\renewcommand{\headrulewidth}{0pt}
\fancyfoot[L]{\sffamily\footnotesize\color{dempet}mattelab.no · Stigningstallfunksjonen · elevark uten fasit}
\fancyfoot[R]{\sffamily\footnotesize\color{dempet}Side \thepage}

\newcommand{\fase}[2]{\tcbox[on line, colframe=#2, colback=white, coltext=#2, boxrule=0.6pt, arc=2pt,
  boxsep=0pt, left=3pt, right=3pt, top=1.5pt, bottom=1.5pt, fontupper=\sffamily\bfseries\footnotesize]{#1}}
\newcommand{\steg}[4]{%
  \par\noindent{\color{rute}\rule{\linewidth}{0.4pt}}\par\nopagebreak
  \noindent\begin{minipage}[t]{2.1cm}\vspace{0pt}\fase{#1}{#2}\end{minipage}%
  \begin{minipage}[t]{\dimexpr\linewidth-2.1cm\relax}\vspace{0pt}#3\end{minipage}\par
  \vspace{#4}}
\newcommand{\del}[4]{%
  \par\noindent{\color{rute}\rule{\linewidth}{0.4pt}}\par\nopagebreak
  \noindent\begin{minipage}[t]{1.1cm}\vspace{0pt}\fase{#1}{#2}\end{minipage}%
  \begin{minipage}[t]{\dimexpr\linewidth-1.1cm\relax}\vspace{0pt}#3\end{minipage}\par
  \vspace{#4}}
\newcommand{\oppgave}[2]{\par\Needspace{12\baselineskip}\vspace{0.6em}%
  {\fontsize{14}{17}\selectfont\bfseries{\color{m1}#1}\enspace #2}\par\vspace{0.2em}}

\begin{document}
\begin{minipage}[t]{0.72\linewidth}\vspace{0pt}
{\sffamily\small\color{dempet}Elevark · R1 · mattelab.no}\par\vspace{0.3em}
{\fontsize{24}{28}\selectfont\bfseries Stigningstallfunksjonen}\par\vspace{1em}
{\sffamily\small Navn: \rule{4.6cm}{0.4pt}\quad Klasse: \rule{1.5cm}{0.4pt}\quad Dato: \rule{1.9cm}{0.4pt}}
\end{minipage}\hfill
\begin{minipage}[t]{0.22\linewidth}\vspace{0pt}\raggedleft
\qrcode[height=2.3cm]{https://mattelab.no/utforskere/stigningstallfunksjonen.html}\par\vspace{0.2em}
{\sffamily\scriptsize\color{dempet}Utforskeren}
\end{minipage}

\vspace{1em}
{\sffamily\small\color{dempet}Utforskeren ligger på \url{https://mattelab.no/utforskere/stigningstallfunksjonen.html}}\par
Hver oppgave har tre steg. \textbf{Gjett} først, uten å bruke utforskeren. \textbf{Test} så i utforskeren. \textbf{Forklar} til slutt med egne ord hvorfor det blir slik. La fasiten være av til du har gjettet.

\oppgave{1}{Skisser grafen til $f'$}
\steg{Gjett}{m1}{Se på grafen til $f(x) = x^3 - 3x$. Skisser grafen til $f'$ på papir. Hvor er $f'$ positiv, negativ og null?}{2.2cm}
\steg{Test}{m2}{Velg $f(x) = x^3 - 3x$ med $x_0 = -2$. Dra $x_0$ sakte mot høyre, eller trykk «Flytt tangenten over hele grafen». Sammenlign sporet med skissen din.}{2.6cm}
\steg{Forklar}{blekk}{Hvorfor er $f'$ null akkurat der $f$ har toppunkt og bunnpunkt? Hvorfor er $f'$ negativ mellom dem?}{3.4cm}

\oppgave{2}{Tre vannrette tangenter}
\steg{Gjett}{m1}{Grafen til $f(x) = \frac{1}{4}x^4 - 2x^2$ ser ut som en W. Hvor mange nullpunkter har $f'$? Hvilken type graf tror du $f'$ har?}{2.2cm}
\steg{Test}{m2}{Velg $f(x) = \frac{1}{4}x^4 - 2x^2$ med $x_0 = -3$. Lag sporet, og finn $x$-verdiene der $f'(x) = 0$.}{2.6cm}
\steg{Forklar}{blekk}{Bruk fortegnet til $f'$ på hver side av nullpunktene til å avgjøre om $f$ har toppunkt eller bunnpunkt der.}{3.4cm}

\oppgave{3}{En funksjon som er sin egen deriverte}
\steg{Gjett}{m1}{Grafen til $f(x) = e^x$ blir brattere og brattere. Hvordan tror du grafen til $f'$ ser ut?}{2.2cm}
\steg{Test}{m2}{Velg $f(x) = e^x$ med $x_0 = -1$. Sammenlign $f(x_0)$ og $f'(x_0)$ for $x_0 = -1$, $0$, $1$ og $2$.}{2.6cm}
\steg{Forklar}{blekk}{Hva er spesielt med $e^x$? Hva er likningen til tangenten i $x = 0$?}{3.4cm}

\oppgave{4}{Logaritmefunksjonen}
\steg{Gjett}{m1}{Grafen til $f(x) = \ln x$ vokser, men blir slakere og slakere. Hva skjer med $f'(x)$ når $x$ blir stor? Og når $x$ nærmer seg $0$?}{2.2cm}
\steg{Test}{m2}{Velg $f(x) = \ln x$ med $x_0 = 0{,}5$. Les av $f'(x_0)$ for $x_0 = 0{,}5$, $1$, $2$ og $4$.}{2.6cm}
\steg{Forklar}{blekk}{Finn et mønster i tallene. Hvilken funksjon tror du $f'$ er?}{3.4cm}

\oppgave{5}{Hull i grafen til $f'$}
\steg{Gjett}{m1}{Hvordan ser grafen til $f'$ ut når $f(x) = |x|$? Hva skjer i $x = 0$?}{2.2cm}
\steg{Test}{m2}{Velg $f(x) = |x|$ med $x_0 = -3$, og lag sporet. Sett så $x_0 = 0$.}{2.6cm}
\steg{Forklar}{blekk}{Hvorfor finnes ikke $f'(0)$, selv om $f(0) = 0$ finnes?}{3.4cm}

\end{document}
